\section{Track G: Recent advances in optimization under uncertainty}\newpage

\begin{talk}
  {Stability of Expected Utility in Bayesian Optimal Experimental Design}% [1] talk title
  {Tapio Helin}% [2] speaker name
  {LUT University}% [3] affiliations
  {tapio.helin@lut.fi}% [4] email
  {}%{Names of coauthors go here, no affiliations of coauthors please, all affiliations will be included in an appendix of 
  %the program book}% [5] coauthors
  {Recent advances in optimization under uncertainty}% [6] special session. Leave this field empty for contributed talks. 
    
   
We explore the stability properties of utility functions in Bayesian optimal experimental design, introducing novel utilities inspired by the analysis. We establish a general framework to study the behavior of expected utility under perturbations in infinite-dimensional setting and prove its convergence properties. Our results provide theoretical guarantees for the robustness of Bayesian design criteria and offer insights into their practical applicability in complex experimental settings.

\medskip

%If you would like to include references, please do so by creating a simple list numbered by [1], [2], [3], \ldots. See example below.
%Please do not use the \texttt{bibliography} environment or \texttt{bibtex} files.
%APA reference style is recommended.
%\begin{enumerate}
% \item[{[1]}] Niederreiter, Harald (1992). {\it Random number generation and quasi-Monte Carlo methods}. Society for Industrial and Applied Mathematics (SIAM).
% \item[{[2]}] Roberts, Gareth O, \& Rosenthal, Jeffrey S. (2002).  Optimal scaling for various Metropolis-Hastings algorithms, \textbf{16}(4), 351--367.
%\end{enumerate}
%
%Equations may be used if they are referenced. Please note that the equation numbers may be different (but will be cross-referenced correctly) in the final program book.
\end{talk}

\begin{talk}
  {Subspace accelerated measure transport methods for fast and scalable sequential
experimental design}% [1] talk title
  {Karina Koval}% [2] speaker name
  {Heidelberg University}% [3] affiliations
  {karina.koval@iwr.uni-heidelberg.de}% [4] email
  {Tiangang Cui, Roland Herzog, Robert Scheichl}% [5] coauthors
  {Recent advances in optimization under uncertainty}% [6] special session. Leave this field empty for contributed talks. 
    
   
In this talk, we present a novel approach for sequential optimal experimental design (sOED) for Bayesian inverse problems governed by expensive forward models with large-dimensional unknown parameters. Our goal is to design experiments that maximize the expected information gain (EIG) from prior to posterior, a task that is computationally challenging. This task becomes more complex in sOED, where we must approximate the incremental expected information gain (iEIG) multiple times in distinct stages, often dealing with intractable prior and posterior distributions. To address this, we propose a derivative-based upper bound for iEIG that guides experimental design and enables parameter dimension reduction through likelihood-informed subspaces. By combining this approach with transport map surrogates for the sequence of posteriors, we develop a unified framework for parameter dimension-reduced sOED. We demonstrate the effectiveness of our approach with numerical examples inspired by groundwater flow and photoacoustic imaging.
\end{talk}

\begin{talk}
  {Randomized quasi-Monte Carlo methods for risk-averse stochastic optimization}% [1] talk title
  {Johannes Milz}% [2] speaker name
  {H. Milton Stewart School of Industrial and Systems Engineering, Georgia Institute of Technology}% [3] affiliations
  {johannes.milz@isye.gatech.edu}% [4] email
  {Olena Melnikov}% [5] coauthors
  {Recent advances in optimization under uncertainty}% [6] special session. Leave this field empty for contributed talks.
    

We establish epigraphical and uniform laws of large numbers for sample-based approximations of law invariant risk functionals. These sample-based approximation schemes include Monte Carlo (MC) and certain randomized quasi-Monte Carlo integration (RQMC) methods, such as scrambled net integration. Our results can be applied to the approximation of risk-averse stochastic programs and risk-averse stochastic variational inequalities. Our numerical simulations empirically demonstrate that RQMC approaches based on scrambled Sobol' sequences can yield smaller bias and root mean square error than MC methods for risk-averse optimization.




\end{talk}

\begin{talk}
  {Efficient expected information gain estimators based on the randomized quasi-Monte Carlo method}% [1] talk title
  {Arved Bartuska}% [2] speaker name
  {King Abdullah University of Science and Technology/RWTH Aachen University}% [3] affiliations
  {arved.bartuska@kaust.edu.sa}% [4] email
  {Andr\'{e} Gustavo Carlon, Luis Espath, Sebastian Krumscheid, Ra\'{u}l Tempone}% [5] coauthors
  
    
   
Efficient estimation of the expected information gain (EIG) of an experiment allows for design optimization in a Bayesian setting. This task faces computational challenges, particularly when the experiment model requires numerical discretization schemes. We demonstrate various methods to make such estimations feasible, combining quasi-Monte Carlo (QMC), randomized QMC (rQMC), and multilevel methods.

Analytical error bounds are made possible by Owen's [1] and He et al.'s [2] work on singular integrands combined with a truncation scheme of the observation noise present in experiment models. Applications from Bayesian experimental design demonstrate the improved convergence behavior of the proposed methods compared to traditional Monte Carlo-based estimators.

\medskip

\begin{enumerate}
 \item[{[1]}] Owen, Art B. (2006). {\it Halton sequences avoid the origin}. SIAM Review, 48:487--503.
 \item[{[2]}] He, Zhijian, \& Zheng, Zhan, \& Wang, Xiaoqun. (2023).{\it On the error rate of importance sampling with randomized quasi-Monte Carlo}. SIAM Journal
on Numerical Analysis,  61(2):10.1137/22M1510121.
\end{enumerate}

\end{talk}
